\textbf{Proposition 3.7} Let $G \subseteq \mathbb{Z}$ be a subgroup. Then $G = d\mathbb{Z}$ for some $d \ge 0$.
\textbf{Proof:}If $G = \{00\}$, then $G = 0\mathbb{Z}$. If not , note that $G$ must contain positive integers: indeed, if $a \in G$ and $a < 0$, then $-a \in G$ and $-a > 0$. We can then let $d$ be the smallest positvie integers in $G$, and I claim $G = d\mathbb{Z}$. The inclusion $\mathbb{Z}\subseteq G$ is clear as G is a subgroup. To verify the inclusion $G \subseteq d\mathbb{Z}$, let $m\in G$ and apply 'division with remainder' to write $$m = dq + r , 0 \le r < d.$$
Since $\in G$ and $d\mathbb{Z}\subseteq G$ and since $G$ is a subgroup, we see that $r = m - dq \in \mathbb{G}$. But d is the smallest positve integer in G, and $\in G$ is smaller than $d$, os $r$ cannot be positve. This shows $r = 0$, that is $m = qd \in d\mathbb{Z} : G\subseteq d\mathbb{Z}$ follows. Thus $G = d\mathbb{Z}$.
The quotient homomorphism $\varPi_n : \mathbb{Z} \to \mathbb{Z}_n$ allows us to establish the analogus result for finite cyclic groups:

\textbf{Propositive 3.8:}
 Let $ n>0 $ be the quotient map, and let $G \subseteq Z_n$ be a subgroup. Then $G$ is the cyclic subgroup of $Z_n$ generated by $d_n$, for some divisor of $d$ of $n$.
\textbf{Proof:}
Let $ \varpi_n: \mathbb{Z} \to \mathbb{Z_n}$ be the quotient map, and consider $G^\prime := \varPi_n^{-1}(G)$.
 By Lemma 3.4, $G^\prime$ is a subgroup of $\mathbb{Z}$ In addition, by Proposition 3.7. $G^\prime$ is a cyclic subgroup of $\mathbb{Z}$
 generated by a non-negative integer $d$. If follows that
 $$ G = \varPi_n(G^\prime) = \varPi(\langle d \rangle) = \langle [d]_n \rangle$$;
thus $G$ is indeed a cyclci subgroup of $Z_n$ generated by a class $d_n$. Furthermore,
since $n \in G^\prime$(because $\varPi_n(n) = [n]_n = [0]_n \in G$) and $G^\prime = d\mathbb{Z}$
we see that $d$ divides $n$ and clamimed . As a consequence of Proposition 6.11, there is a bijection between
the set of subgroups of $\mathbb{Z}_n$ and the set of positive divisors of n.
For example, $\mathbb{Z}_12$ has exactly 6 subgroups, because 12 has 6 positive divisors: $1,2,3,4,6 and 12.$
 Here is the corresponding list of subgroups: \\
$\langle [1]_12\rangle = \{[0]_{12}, [1]_{12}, [2]_{12}, [3]_{12}, [4]_{12}, [5]_{12},[6]_{12} ,[7]_{12},[8]_{12}, [9]_{12}, [10]_{12}, [11]_{12}\}$ \\
$\langle [2]_12\rangle = \{[0]_{12}, [2]_{12},[4]_{12},[6]_{12} ,[8]_{12}, [10]_{12}\}$ \\
$\langle [3]_12\rangle = \{[0]_{12}, [3]_{12}, [6]_{12}, [9]_{12}\}$ \\
$\langle [4]_12\rangle = \{[0]_{12}, [4]_{12}, [8]_{12},\}$ \\
$\langle [6]_12\rangle = \{[0]_{12}, [6]_{12}\}$ \\
$\langle [12]_12\rangle = \{[0]_{12}\}$ \\

Also note that if $d_1,d_2$ are both divisors of $n$, and $d_1 | d_2$ then $\langle[d]_n\rangle \subseteq \langle [d_1]_n\rangle$. That is , the correspondence between subgroups of $\mathbb{Z}_n$ and divisors of n preserves the natural lattice structure carried by these sets. We can draw these lattices for $\mathbb{Z}_6$ as follows
\\
\begin{center}
\begin{tikzpicture}[
    % Define a style for the small empty circles at each vertex
    dot/.style={circle, draw=black, inner sep=0pt, minimum size=4pt, fill=white},
    scale=1.2 % Adjust the overall size if needed
]

% ==============================
% Left Diagram
% ==============================
\begin{scope}[xshift=0cm]
    % Nodes
    \node[dot, label=above:{$1$}] (n1) at (0, 2) {};
    \node[dot, label=left:{$2$}]  (n2) at (-1.2, 0) {};
    \node[dot, label=right:{$3$}] (n3) at (1.2, 0) {};
    \node[dot, label=below:{$6$}] (n6) at (0, -2) {};

    % Edges
    \draw (n1) -- (n2);
    \draw (n1) -- (n3);
    \draw (n1) -- (n6); % Middle vertical line
    \draw (n2) -- (n6);
    \draw (n3) -- (n6);
\end{scope}

% ==============================
% Right Diagram
% ==============================
\begin{scope}[xshift=6.5cm]
    % Nodes
    \node[dot, label=above:{$\{[0], [1], [2], [3], [4], [5]\}$}] (top) at (0, 2) {};
    \node[dot, label=left:{$\{[0], [2], [4]\}$}]                 (left) at (-1.5, 0) {};
    \node[dot, label=right:{$\{[0], [3]\}$}]                     (right) at (1.5, 0) {};
    \node[dot, label=below:{$\{[0]\}$}]                          (bottom) at (0, -2) {};

    % Edges (Note: No middle vertical line in this diagram)
    \draw (top) -- (left);
    \draw (top) -- (right);
    \draw (left) -- (bottom);
    \draw (right) -- (bottom);
\end{scope}

\end{tikzpicture}
\end{center}
\section{Monomporphisms}
Let $G$ and $G^\prime$ be gourps and let $\varphi : G \to G^\prime$ be a group homomorphism is said to be a monomorphism if it is an injective group homomorphism. The following characterics monomorphisms.

\textbf{Poroposition 3.9}
The following are equivalent
\begin{enumerate}
    \item $\varphi$ is a monomorphism
    \item $Ker\varphi = \{e_G\}$
    \item $\varphi : G \to G^\prime$ is injective (as a set-function)
\end{enumerate}

\text{Proof:} $(a) \implies (b)$: Assume (a) holds, and consider the two parallet composition,
\[
\begin{tikzcd}[column sep=large]
\ker \varphi \arrow[r, shift left=0.7ex, "i"] \arrow[r, shift right=0.7ex, "e"'] & G \arrow[r, "\varphi"] & G',
\end{tikzcd}
\]
where $i$ is the inclusion and $e$ is the trivial map.
Both $\varphi \circ i$ and $\varphi \circ e$ are the
trivial maps since $\varphi$ is a monomorphism, this
implies $i = e$. But $i=e$ implies that $ker\varphi = \{e_G\}$,
so that $(b)$ holds.\\
$(b) \implies (c):$ Assume $ker\varphi = \{e_G\}$. Then
$$\varphi(g_1) = \varphi(g_2) \implies \varphi(g_1)\varphi(g_2)^{-1} = e_G^\prime \implies \varphi(g_1g_2^{-1} = e_G^\prime)$$
$$\implies g_1g_2^{-1} \in ker\varphi \implies g_1g_2^{-1} = e_G \implies g_1 = g_2$$
This shows that $\varphi$ is injective as required.
$(c) \implies (a):$ If $\varphi$ is injective, then it satisfies the defining
property for monomorphism as a set function:that is, for any set $X$ and any two set-functions $\alpha^\prime, \alpha^{\prime \prime}: X \to G$,
$$\varphi \circ \alpha^\prime = \varphi \circ \varphi^{\prime \prime} \iff \alpha^\prime = \alpha^{\prime \prime}$$
This must hold in particular if $X$ has a group structure and $\alpha^\prime$ and $\alpha^{\prime \prime}$ are group homomorphism, so $\varphi$ is a group monomorphism.

\section{Quotient Groups}
\textbf{Normal Subgroups:}
\textbf{Definition 4.1} A subgroup $N$ of a group $G$ is normal if $\forall g \in G, \forall n \in N, gng^{-1} \in N.$
Note that every subgroup of a commutative group is normal because
$\forall g \in G, gng^{-1} = n \in N$. However, in general not all subgroups are normal examples may be found already in $S_3$
There exist noncommutative group in which every subgroup is normal(one exampe is the 'quaternionic group)' $Q_8$.
\textbf{4.2}: If $\varphi : G \to G^\prime$ is any group homomorphism, then $ker\varphi$ is a normal subgroup of $G$.
\textbf{Proof:}
As we already know, $ker \varphi$ is a subgroup of $G$; to verify it is normal, note that $\forall g \in G, \forall n \in ker\varphi$
$$\varphi(gng^{-1}) = \varphi(g)\varphi(n)\varphi(g)^{-1} = \varphi(g)e_G\varphi(g)^{-1} = e_G\prime$$. Proving that $gng^{-1} \in ker\varphi$. So in shork kernel implies normal.
Convenienty we can express normality as follows: if $g\in G$ and $A \subseteq G $ is any subset , we denote by $gA,Ag$ resprectively, the following subsets of G:
$$gA := \{h\in G | (\exists a \in A): h = ga\}$$ $$Ag := \{h\in G | (\exists a \in A): h = ag\}$$
Then the normality condition can be expressed by $$(\forall g \in G): gNg-1 \subseteq N$$, or in anumber of other ways:
$$gNg^{-1} = N or  gN \subseteq Ng or gN = Ng$$ for all $g\in G$. Note that all of there are equivalent conditions. Keep in mind that $gN = Ng$ does not mean that g commutes with every element of N ; it means that if $n \in N$, then there are elements $n',n'' \in N$, in general different from n, such that $gn' = n'g$. (So that $gN \subseteq Ng$) and $ng = gn''$ (so that $Ng \subseteq gN$).

\section{Quatient group}
We consider an equivalence relation $\sim$ on (the set underlying) a gourp $G$, we seek a group $G/\sim$ and a group homomorphism $\varPi : G \to G/\sim.$
It is natural to try to construct the group $G/\sim$ by defining an operation $\cdot$ on the set $G/\sim$. This construction is constrained by the requirement that the quotient map $\varPi: G \to G/\sim$ be a group homomorphism
because if $[a] = \varPi(a), [b] = \varPi(b)$ are elements of $G/\sim$ (that is, equivalence classes with respect to $\sim$), then the homomorphism condition forces
$$[a]\cdot [b] = \varPi(a) \cdot \varPi(b) = \varPi(ab) = [ab]$$. However we need to ask ourselves whether the operationis well defined(In the first factor). This leads to a conditins on the equivalence relation which we will discuss. For the operation to be well-defined, it is necessary that if $[a] = [a']$, then $[ab] = [a'b]$ regardless of what b is ; that is
$$(\forall g \in G): a \sim a' \implies ag \sim a'g$$. Similaly for the operation to be well defined in the second factor we need
$$(\forall g \in G): a \sim a' \implies ga = ga'$$

\textbf{Proposition 4.3} With the above notation, the operation $[a]\cdot[b]:= [ab]$
defines a group structure on $G/\sim$ if and only if $\forall a, a', g \in G$
$$a \sim a' \implies ga \sim ga' \quad and \quad ag \sim a'g.$$

\underline{Proof:}

We have already noted that the stated condition is necessary. To prove it is sufficient, with the stated congruences, assume $\forall a, a', g \in G$.
$$a \sim a' \implies ga \sim ga' \text{ and } ag \sim a'g$$

Then the operation
$$[a] \cdot [b] := [ab]$$
is well-defined and we have to verify that it defines a group structure on $G/\sim$. The associativity of $\cdot$ is inherited from the associativity of $G$: $\forall a,b,c \in G$
\begin{align*}
([a] \cdot [b]) \cdot [c] & = [ab] \cdot [c] = [(ab)c] = [a(bc)] \\
&= [a] \cdot [bc] \\
&= [a] \cdot ([b] \cdot [c]).
\end{align*}

The class $[e_G]$ is an identity with respect to this operation: $\forall g \in G$
$$[g] \cdot [e_G] = [ge_G] = [g], \quad [e_G] \cdot [g] = [e_Gg] = [g].$$

The class $[g^{-1}]$ is the inverse of $[g]$:
$$[g^{-1}] \cdot [g] = [g^{-1}g] = [e_G], \quad [g] \cdot [g^{-1}] = [gg^{-1}] = [e_G].$$

Thus $G/\sim$ is indeed a group.